\section{Evidence Hierarchy}
\label{app:evidence_hierarchy}

The evidence hierarchy table is presented in the main text as Table~\ref{tab:evidence_hierarchy_main}. The categories, supporting evidence, and limitations are defined there.

\ArtifactTable{../artifacts/tables/mnl_sensitivity_summary.tex}{MNL-regime sensitivity summary separating current evidence from completed sensitivity reruns.}

\ArtifactTable{../artifacts/tables/filter_validity_summary.tex}{Filter-validity diagnostics separating displayed-offer and selected-offer ETA/IVT errors.}

\section{Lambert-W Reference Transform (Supplementary)}
\label{app:lambertw}

We summarize the revenue-only common-offset reference problem used to motivate the implemented Lambert-W-inspired heuristic. This appendix is not a proof of optimality for the full simulator objective, which includes heterogeneous route costs, clipped prices, route-aware menu selection, and an outside option.

\paragraph{Setup.} Let $\mathcal{M}$ be the displayed menu. The operator sets a common price offset $\delta$ applied to all bundles, so that the price of bundle $b$ is $p_b = \hat{c}(b) - R - \delta$, where $R$ is the revenue target and $\hat{c}(b)$ is the predicted cost. The utility of bundle $b$ at offset $\delta$ is
\[
  u_b(\delta) = u_b^0 + \beta\,(\hat{c}(b) - R - \delta),
\]
where $u_b^0$ is the utility evaluated at zero price and $\beta < 0$ is the price sensitivity. With the outside option normalized to $u_0 = 0$, the MNL choice probability is
\[
  \pi_b(\delta) = \frac{e^{u_b(\delta)}}{e^{u_0} + \sum_{b' \in \mathcal{M}} e^{u_{b'}(\delta)}},
\]

\paragraph{Reference objective.} Following the revenue-management setting of \citet{talluri2004revenue}, consider the expected displayed revenue under the shared offset,
\[
  \mathcal{R}(\delta) = \sum_{b \in \mathcal{M}} \pi_b(\delta)\,p_b(\delta).
\]
The outside option remains in the denominator through $e^{u_0}$, so the total probability assigned to displayed bundles is generally below one. The full simulator objective used in the experiments is instead an expected net-profit surrogate,
\[
  \widehat{\Pi}(\mathcal{M}) = \sum_{b \in \mathcal{M}} \pi_b(\mathcal{M})\{p_b-\hat c_{\mathrm{sys}}(b)\},
\]
as defined in Eq.~\ref{eq:menu_value}. We do not simplify that objective to a linear expression in $\delta$; doing so would incorrectly remove the outside-option denominator and would also ignore the route-cost heterogeneity that the simulator evaluates.

\paragraph{Lambert-W closed form.} When the objective is the expected revenue
$\sum_b \pi_b(\delta)\,p_b$ (equivalently, when costs are treated as fixed), the
first-order condition with respect to $\delta$ yields a transcendental equation
solvable via the Lambert W function. This revenue-based derivation motivates the implemented common-offset heuristic, but it should not be read as a proof that the clipped pricing rule is globally profit-optimal under the full evaluated objective. Define
\[
  S = e^{u_0} + \sum_{b \in \mathcal{M}} e^{u_b^0 + \beta(\hat{c}(b) - R)}.
\]
Under the shared-offset parameterization, the first-order condition can be rearranged into the standard Lambert-W form
\[
  w e^w = S/e,
\]
with solution $w = W_0(S/e)$ on the principal branch. Mapping back to the price offset gives
\begin{equation}
  \delta^* = \frac{W_0(S/e) + 1}{\beta}.
  \tag{A.1}\label{eq:lambertw_deriv}
\end{equation}
The price for each bundle is then clipped to the platform bounds:
\[
  p_b = \operatorname{clip}\!\bigl(\hat{c}(b) - R - \delta^*,\; p_{\min},\; p_{\max}\bigr).
\]
This appendix records the closed-form motivation for a narrow reference problem only. The implemented policy in the main text adds route-aware menu selection, heterogeneous predicted costs, price clipping, and downstream HGS re-optimization, so Eq.~\eqref{eq:lambertw_deriv} should be interpreted as a heuristic reference transform rather than an exact optimality certificate for the final simulator objective.

\paragraph{Conditions for validity.} The common-offset assumption requires that all
passengers share the same price sensitivity $\beta$ and that utility is additively
separable in price. Both hold by construction in the MNL simulator
(Table~\ref{tab:mnl_params}). The Lambert-W solution is real-valued whenever
$S > 0$, which is guaranteed as long as at least one bundle has finite utility. The
home bundle's large positive utility term ($u^{\text{home}} = 3.2$) ensures
$S \gg 0$ in all practical cases, which is why removing the home bundle from the
pricing set causes $S$ to collapse and $\delta^*$ to fall below $p_{\min}$,
collapsing all OOH prices to the floor.

\section{Scoring-Weight Sensitivity Analysis}
\label{app:sensitivity}

Table~\ref{tab:sensitivity} reports an \emph{archival} single-split diagnostic from an earlier pipeline configuration: the net-profit gap of the no-filter heuristic versus full display when each of the seven scoring and menu parameters is individually perturbed by $\pm 50\%$ from its baseline value on the \texttt{rc\_0\_to\_1} split. These numbers are retained only to show relative directional sensitivity within that historical setting; they are not directly comparable to the current outside-option results in Section~\ref{sec:results}.

\begin{table}[h]
\centering
\caption{Archival single-split sensitivity diagnostic from an earlier pipeline configuration. Baseline gap in that setting is $+659.6$; the table is included only to show relative directional sensitivity within that historical setup, not to support the main claims of the paper.}
\label{tab:sensitivity}
\begin{tabular}{llrrr}
\hline
Parameter & Description & $-50\%$ & Baseline & $+50\%$ \\
\hline
$\mu_m$ & Margin weight & 659.6 & 657.7 & 663.2 \\
$\mu_w$ & Walking-distance weight & 663.2 & 657.9 & 661.5 \\
$\mu_t$ & Pickup-time deviation weight & 656.0 & 657.7 & 659.4 \\
$\mu_v$ & In-vehicle time weight & 659.6 & 659.6 & 659.6 \\
$\mu_r$ & Route-delay penalty weight & 578.0 & 659.6 & 657.5 \\
$\mu_q$ & Capacity-risk penalty weight & 597.5 & 659.6 & 645.6 \\
$k$ & Menu size cap ($k=2, 3, 5$) & 607.8 & 659.6 & 637.8 \\
\hline
\end{tabular}
\end{table}

Within that historical setting, the route-delay penalty weight $\mu_r$ and the capacity-risk penalty weight $\mu_q$ exhibit the largest downward sensitivity when reduced to 50\% of their baseline values. The IVT weight $\mu_v$ is insensitive across the tested range, reflecting the low variance in predicted in-vehicle times for the RC instance. The menu size $k$ shows a modest negative effect at $k=5$, likely because larger menus include marginal bundles that dilute the choice probability of the highest-profit options. Re-running this sweep under the current evaluation pipeline remains future work.

\section{Predictor Validation}
\label{app:predictor_validation}

Table~\ref{tab:predictor_mae} reports the mean absolute error (MAE) of the CNN-2d predictor and a naive mean-prediction baseline for each output head, evaluated on a held-out 20\% partition of the training buffer across all six RC splits.

\begin{table}[h]
\centering
\caption{Predictor MAE on held-out buffer data (mean $\pm$ std across 6 RC splits). CNN-2d reduces cost MAE by 15.3\% relative to the naive mean baseline. ETA and IVT CNN MAE is higher than the naive baseline because the CNN is trained on noisy online data; in deployment, ETA and IVT predictions are blended with heuristic estimates (blend weight 0.35 CNN, 0.65 heuristic).}
\label{tab:predictor_mae}
\begin{tabular}{lrrr}
\hline
Predictor & Cost MAE & ETA MAE & IVT MAE \\
          & (monetary units) & (seconds) & (seconds) \\
\hline
CNN-2d     & $5.00 \pm 0.59$ & $5{,}703 \pm 286$ & $5{,}275 \pm 247$ \\
Naive mean & $5.91 \pm 0.57$ & $1{,}863 \pm 66$  & $1{,}860 \pm 67$  \\
\hline
\end{tabular}
\end{table}

The naive mean baseline predicts the training-set mean for each head; it serves as the simplest possible zero-information predictor. The CNN-2d reduces cost prediction error by 15.3\% relative to this baseline. The cost head is the primary input to the system-evaluation and menu-value calculations in Eqs.~\ref{eq:syscost} and \ref{eq:menu_value}, so improvement in cost prediction directly benefits menu ranking quality.

For ETA and IVT, the CNN MAE is higher than the naive baseline because the predictor is trained online on data collected during early training episodes, when the routing state is highly variable. In deployment, ETA and IVT predictions are blended with heuristic estimates derived from cheapest-insertion cost computations (blend weight 0.35 CNN, 0.65 heuristic; parameter \texttt{menu\_time\_prediction\_blend}), which substantially reduces the effective prediction error applied to the score function.

\section{Split-Level Uncertainty and Policy-Conditioned Diagnostics}
\label{app:split_uncertainty}

Table~\ref{tab:policy_gap_uncertainty_summary} reports the split-level uncertainty summary used in the main text. The sampling unit throughout is the split-level paired mean net-profit gap relative to full display. RC therefore supports modest split-level uncertainty reporting with $n=6$ split pairs, whereas Austin and Seattle remain descriptive because each city contributes only $n=2$ split pairs.

Table~\ref{tab:rc_split_gap_summary} lists the six RC split-level paired gaps together with their within-split episode-level CI half-widths. These within-split intervals are useful for diagnosing heterogeneity across the replayed traces inside a fixed split, but the paper's uncertainty statements are based on the split-level paired means rather than on pooling episodes across splits.

\ArtifactTable{../artifacts/tables/rc_main_optout_split_gap_table.tex}{RC outside-option split-level paired gaps against full display. Within-split CI half-width is computed from episode-level pairing inside each split.}

Table~\ref{tab:city_split_gap_summary} gives the corresponding Austin and Seattle split-level paired gaps. With only two split pairs per city, these numbers should be read as descriptive evidence about directional stability rather than as precise interval estimates.

\ArtifactTable{../artifacts/tables/city_split_gap_table.tex}{Austin and Seattle split-level paired gaps against full display. Each city has only two split pairs, so these rows are descriptive.}

Finally, Table~\ref{tab:city_policy_diagnostic_summary} separates policy-conditioned displayed-offer diagnostics from realized passenger-facing outcomes. The displayed ETA/IVT MAE and false-negative pruning columns diagnose how each policy changes the offered-candidate distribution seen by the predictor. Acceptance rate and consumer surplus, by contrast, remain realized evaluation outcomes under replayed traces.

\ArtifactTable{../artifacts/tables/city_policy_diagnostic_summary.tex}{City-level policy-conditioned displayed-offer diagnostics and realized passenger-facing outcomes.}

\section{Computational Cost}
\label{app:computational_cost}

Table~\ref{tab:computational_cost} summarises the wall-clock computational cost of the three main phases: predictor training, menu-construction inference, and HGS back-end routing.

\begin{table}[h]
\centering
\caption{Computational cost summary for the RC instance (90 customers, 10 meeting points, 15 vehicles). Training and inference times are measured on a standard CPU (no GPU used). HGS time limits are upper bounds configured per episode; actual routing time is typically below the limit.}
\label{tab:computational_cost}
\begin{tabular}{llr}
\hline
Phase & Metric & Value \\
\hline
Predictor training & Training episodes & 12 \\
                   & Steps per episode & ${\approx}700$ \\
                   & Wall-clock per episode & ${\approx}26.5$ s \\
                   & Total training time & ${\approx}318$ s ($\approx 5.3$ min) \\
                   & Initial phase epochs & 12 \\
                   & Initial phase total time & ${\approx}0.4$ s \\
\hline
Menu-construction inference & Mean time per request & $44.9$ ms \\
                            & (CNN forward pass + insertion cost) & \\
\hline
HGS routing & Reopt time limit (per step) & $1.1$ s \\
            & Final reopt time limit & $1.5$ s \\
            & Mean step time (menu + HGS) & ${\approx}43$ ms \\
\hline
\end{tabular}
\end{table}

The menu-construction inference time (44.9 ms per request) includes the CNN forward pass and the cheapest-insertion cost computation used to estimate marginal routing cost. HGS routing is invoked periodically (every \texttt{reopt} steps) rather than at every request; the mean step time of $\approx 43$ ms reflects the amortised cost across steps with and without routing calls.

\section{Reproducibility}
\label{app:reproducibility}

The project keeps the experiment logic and the paper layer synchronized through a small set of public commands:

\begin{verbatim}
python scripts/run_study.py --study rc_main_optout
python scripts/run_study.py --study phase22_eta_compare
python scripts/run_study.py --study phase22_uptake_calibration
python scripts/run_study.py --study phase23_pricing_baselines
python scripts/run_study.py --study phase29_exact_greedy_gap
python scripts/run_study.py --study phase30_robust_filtering
python scripts/run_study.py --study phase31_uptake_menu_value
python scripts/run_study.py --study austin_main
python scripts/run_study.py --study seattle_main
python scripts/extract_phase22_results.py
python scripts/extract_phase23_results.py
python scripts/extract_phase24_results.py
python scripts/build_artifacts.py --study rc_main_optout
python scripts/build_artifacts.py --study phase29_exact_greedy_gap
python scripts/build_artifacts.py --study phase30_robust_filtering
python scripts/build_artifacts.py --study phase31_uptake_menu_value
python scripts/build_artifacts.py --study austin_main
python scripts/build_artifacts.py --study seattle_main
python scripts/build_manuscript.py
\end{verbatim}

For local smoke verification, the repository also ships a minimal manifest:

\begin{verbatim}
python scripts/run_study.py --study smoke_rc
\end{verbatim}

The manuscript consumes generated tables and figures from \texttt{artifacts/}, while full raw experiment outputs remain under \texttt{outputs/} and are excluded from version control.

\section{Consumer Surplus Computation}
\label{app:consumer_surplus}

Under the MNL choice model, the expected consumer surplus for a single request $i$ facing menu $\mathcal{M}_i$ is:
\begin{equation}
\text{CS}_i = \frac{1}{\alpha_p} \ln\!\left(\sum_{j \in \mathcal{M}_i} e^{V_{ij}} + e^{V_{i0}}\right)
\end{equation}
where $V_{ij}$ is the deterministic utility of bundle $j$, $V_{i0}$ is the deterministic utility of the outside option, and $\alpha_p$ is the price sensitivity coefficient (the marginal utility of money). When the passenger chooses the outside option, the realized surplus is normalized to zero; when a menu bundle is chosen, the realized surplus equals $\frac{1}{\alpha_p}(\ln e^{V_{ij^*}} - \ln e^{V_{i0}})$ for the chosen bundle $j^*$.

We report two surplus measures:
\begin{itemize}
\item \textbf{All-request surplus:} $\overline{\text{CS}}_{\text{all}} = \frac{1}{N}\sum_{i=1}^{N} \text{CS}_i$, averaged over all $N$ requests including those that chose the outside option. By construction, requests choosing the outside option contribute zero surplus, so this measure reflects both the probability of acceptance and the quality of accepted bundles.
\item \textbf{Accepted-user surplus:} $\overline{\text{CS}}_{\text{acc}} = \frac{1}{N_{\text{acc}}}\sum_{i \in \mathcal{A}} \text{CS}_i$, averaged only over the $N_{\text{acc}}$ accepted requests $\mathcal{A}$ that selected a menu bundle. This measure reflects the average passenger experience conditional on using the service.
\end{itemize}

These two measures can diverge when policies change acceptance rates. A filter that removes low-value bundles may increase $\overline{\text{CS}}_{\text{acc}}$ (higher-quality remaining offers) while decreasing $\overline{\text{CS}}_{\text{all}}$ (fewer accepted passengers). The MNL parameters are fixed design choices rather than estimated from passenger data, so reported surplus values are internally consistent model outputs rather than welfare measurements from revealed preferences.
